\documentclass{article}
\usepackage{graphicx}
\usepackage{amsmath}
\usepackage{geometry} % Required for customizing margins

\geometry{margin=1in} % Sets 1-inch margins on all sides
\title{EMCH 360 Assignment \#3}
\author{Jacob C. Vaught}
\date{October 15th 2023}


\begin{document}

\maketitle
\section*{Problem \#1}

Water flows from a pressurized tank, through a 6-in. diameter pipe, exits from a 2-in. diameter nozzle, and rises 20 ft. above the nozzle. Determine the pressure in the tank if the flow is steady, frictionless, and incompressible.

\subsection*{Solution}

Assume top of tank is at same elevation as nozzle. 
\newline\newline
Bernoulli's Equation:
\[
P_1 + \frac{1}{2} \rho v_1^2 + \rho gh_1 = P_2 + \frac{1}{2} \rho v_2^2 + \rho gh_2
\]
Substituting in the values 20 ft + 2 ft for the water apogee, and 2 ft - 0.1 ft for the tank water surface:

\[
\begin{aligned}
P_{\text{tank}} &= 62.4 \text{lbm/ft}^3 \times 32.2 \text{ft/s}^2 \times \frac{1}{32.2} \times \left( (20+2)  \text{ft} - (2-0.1) \text{ft} \right) \\
P_{\text{tank}} &= 1254.24 \frac{\text{lbf}}{\text{ft}^2} \\
P_{\text{tank}} &= 8.714 \frac{\text{lbf}}{\text{in}^2} \\
\end{aligned}
\]

This will give the gauge pressure inside the tank. To get the absolute pressure, you would need to add the atmospheric pressure to the result.
\newpage
\section*{Problem \# 2}
A fire hose nozzle has a diameter of 1.125 in. According to some fire codes, the nozzle must be capable of delivering at least \( Q = 250 \) gal/min. If the nozzle is attached to a 3-in.-diameter hose, what pressure must be maintained just upstream of the nozzle to deliver this flowrate?

\subsection*{Solution}
To find the pressure just upstream of the nozzle, we can use Bernoulli's equation under the assumption of steady, incompressible flow and no change in height. 

\[
P_1 + \frac{1}{2} \rho v_1^2 = P_2 + \frac{1}{2} \rho v_2^2
\]

Here, \(P_1\) and \(P_2\) are the pressures just upstream and downstream of the nozzle respectively, \(v_1\) and \(v_2\) are the corresponding velocities, and \(\rho\) is the density of water.

First, we can find the velocities \(v_1\) and \(v_2\) by using the flow rate \(Q\) and the given diameters. The flow rate \(Q\) in \( \text{ft}^3/\text{sec} \) can be calculated as follows:

\[
Q = 250 \, \text{gal/min} \times \frac{1 \, \text{ft}^3}{7.481 \, \text{gal}} \times \frac{1 \, \text{min}}{60 \, \text{s}} \approx 0.55696 \, \text{ft}^3/\text{s}
\]
\newline
Using \( Q = A v \), we can find \( v_1 \) and \( v_2 \):

\[
v_1 = \frac{0.55696}{\pi ( \frac{3}{12} )^2 / 4} \approx 11.34 \, \text{ft/s}
\]
\[
v_2 = \frac{0.55696}{\pi ( \frac{1.125}{12} )^2 / 4} \approx 80.685 \, \text{ft/s}
\]
\newline
The density of water is approximately \(62.4 \, \text{lbm/ft}^3\).
\newline\newline
Rearranging Bernoulli's equation to solve for \( P_1 \), we get:

\[
P_1 = P_2 + \frac{1}{2} \rho (v_2^2 - v_1^2)
\]
\newline
Since \( P_2 \) is atmospheric pressure, we can consider it to be zero for this problem.

\[
P_1 = \frac{1}{2} \times 62.4 \, \text{lbm/ft}^3 \times \frac{1}{32.2} \times (80.685^2 - 11.34^2) \, \text{ft}^2/\text{s}^2 \approx 6183.3 \, \text{lbm/ft s}^2 \approx 42.94 \, \text{psi}
\]

\newpage
\section*{Problem \# 3}
The circular stream of water from a faucet is observed to taper from a diameter of \(20 \, \text{mm}\) to \(10 \, \text{mm}\) in a distance of \(50 \, \text{cm}\). Determine the flowrate \(Q\).

\subsection*{Solution}

First, let's find the areas \(A_1\) and \(A_2\) for the two sections of the stream:
\[
A_1 = \frac{\pi (20 \, \text{mm})^2}{4} = 314.16 \, \text{mm}^2
\]
\[
A_2 = \frac{\pi (10 \, \text{mm})^2}{4} = 78.54 \, \text{mm}^2
\]
The flow rate \(Q\) can be expressed in terms of velocities \(V_1\) and \(V_2\) as:
\[
Q = A_1 V_1 = A_2 V_2
\]
From this, we can find the relation between \(V_1\) and \(V_2\):
\[
V_2 = \frac{A_1}{A_2} V_1
\]
We assume \(P_1 = P_2 = 0\) (pressure at both points is atmospheric), and the height difference \(h\) is \(50 \, \text{cm}\). Using Bernoulli's equation:
\[
\frac{V_1^2}{2} = \frac{V_2^2}{2} + gh
\]
Substituting \(V_2\) in terms of \(V_1\):
\[
\frac{V_1^2}{2} = \frac{(A_1/A_2)^2 V_1^2}{2} + gh
\]
Solving for \(V_1\):
\[
V_1 = \sqrt{\frac{2gh}{1-(A_1/A_2)^2}}
\]
\[
V_1 = \sqrt{\frac{2\times9.81\times-0.5}{1-4^2}}
\]
\[
V_1 = 0.8087 \frac{m}{s}
\]
Finally, we can find the flowrate \(Q\):
\[
Q = A_1 V_1 \approx 314.16 \, \text{mm}^2 \times 0.8087 \, \text{m/s} \approx 0.254 \times 10^{-3} \, \text{m}^3/\text{s}
\]





\newpage

\section*{Problem \# 4}
Water exits a pipe as a free jet and flows to a height \( h \) above the exit plane as shown in Fig. P3.71. The flow is steady, incompressible, and frictionless.
\begin{itemize}
    \item The pipe is 4 inches at section 1, horizontally flowing into a vertical pipe which is 6 inches in diameter.
    \item The distance from the middle of the section 1 pipe to the exit plane is 8 ft.
    \item The exit velocity at the exit plane is 16 ft/s.
\end{itemize}

\subsection*{Solution}

\subsubsection*{(a) Determine the height \( h \)}
We can use the energy equation for incompressible, frictionless flow:

\[
h = \frac{V^2}{2g}
\]
Where \( V = 16 \, \text{ft/s} \) and \( g = 32.2 \, \text{ft/s}^2 \).

\[
h = \frac{(16 \, \text{ft/s})^2}{2 \times 32.2 \, \text{ft/s}^2} = \frac{256 \, \text{ft}^2/\text{s}^2}{64.4 \, \text{ft/s}^2} = 3.98 \, \text{ft}
\]

\subsubsection*{(b) Determine the velocity and pressure at section \( 1 \)}
We can use the continuity equation:

\[
A_1 V_1 = A_2 V_2
\]
Where \( A_1 = \pi \left( \frac{4 \, \text{in}}{12 \, \text{ft/in}} \right)^2 / 4 = 0.0873 \, \text{ft}^2 \) and \( A_2 = \pi \left( \frac{6 \, \text{in}}{12 \, \text{ft/in}} \right)^2 / 4 = 0.196 \, \text{ft}^2 \).

\[
V_1 = \frac{A_2 V_2}{A_1} = \frac{0.196 \, \text{ft}^2 \times 16 \, \text{ft/s}}{0.0873 \, \text{ft}^2} = 35.9 \, \text{ft/s}
\]
\newline\newline
We can find the pressure at section 1 using Bernoulli's equation:

\[
P_1 = P_2 + \frac{1}{2} \rho (V_1^2 - V_2^2) - \rho h
\]
Assuming atmospheric pressure at the exit, \( P_2 = 0 \), \( \rho = 1.94 \, \text{slug/ft}^3 \), \( V_1 = 35.9 \, \text{ft/s} \) and \( V_2 = 16 \, \text{ft/s} \).

\[
P_1 = \frac{1}{2} \times 1.94 \, \text{slug/ft}^3 \times ( (35.9 \, \text{ft/s})^2 - (16 \, \text{ft/s})^2 ) - 62.4 \, \text{lbm/ft}^2 \times 8 \, \text{ft}
\]
\[
= 0.97 \, \text{slug/ft}^3 \times (1288.1 - 256) \, \text{ft}^2/\text{s}^2 - 499.2 
\]
\[
= 509.6 \, \text{lbf/ft}^2
\]
\newpage
\section*{Problem \# 5}

Water is siphoned from a large tank and discharges into the atmosphere through a 2-in.-diameter tube. The end of the tube is 3 ft below the tank bottom, and viscous effects are negligible. (a) Determine the volume flowrate from the tank. (b) Determine the maximum height, \(H\), over which the water can be siphoned without cavitation occurring. Atmospheric pressure is \(14.7 \, \text{psia}\), and the water vapor pressure is \(0.26 \, \text{psia}\). The height of the tank is \(9 \, \text{ft}\).

\subsection*{Solution}

First, let's denote the atmospheric pressure as \(P_{\text{atm}} = 14.7 \, \text{psia}\), water vapor pressure as \(P_{\text{vap}} = 0.26 \, \text{psia}\), the gravitational constant as \(g = 32.2 \, \text{ft/s}^2\), and the density of water as \(\rho = 62.4 \, \text{lbm/ft}^3\).

\subsubsection*{(a) Volume Flowrate}

Applying Bernoulli's equation between the free surface of the tank and the tube outlet, we get

\[
P_1 + \frac{1}{2} \rho V_1^2 + \rho g h_1 = P_2 + \frac{1}{2} \rho V_2^2 + \rho g h_2
\]
Since the velocity \(V_1\) in the tank is much smaller than \(V_2\) in the tube, \(V_1\) can be ignored. The pressures \(P_1\) and \(P_2\) are both atmospheric pressures.

\[
V_2 = \sqrt{2g(h_1 - h_2)}
\]
Substituting \(h_1 = 9 \, \text{ft}\) and \(h_2 = -3 \, \text{ft}\),

\[
V_2 = \sqrt{2 \times 32.2 \, \text{ft/s}^2 \times (9 \, \text{ft} - (-3 \, \text{ft}))} = \sqrt{2 \times 32.2 \times 12} \approx 27.8 \, \text{ft/s}
\]
The flow area \(A_2\) is \( \pi \left(\frac{d}{2}\right)^2 \).

\[
A_2 = \pi \left( \frac{2 \, \text{in}}{2} \times \frac{1 \, \text{ft}}{12 \, \text{in}} \right)^2 = \frac{\pi}{144} \, \text{ft}^2
\]
Thus, the volume flow rate \(Q\) is \( A_2 \times V_2 \).

\[
Q = \frac{\pi}{144} \times 27.8 \approx 0.606 \, \text{ft}^3/\text{s}
\]

\subsubsection*{(b) Maximum Height \(H\)}

For cavitation to be avoided, the pressure at the top of the siphon must be greater than the vapor pressure.

\[
P_1 - \rho g h > P_{\text{vap}}
\]

\[
H < \frac{P_1 - P_{\text{vap}}}{\rho g} = \frac{(14.7 - 0.256) \, \text{psia}}{62.4 \, \text{lbm/ft}^3 \times 32.2 \, \text{ft/s}^2} \times 144 \, \text{in}^2/\text{ft}^2 \approx 33.3 \, \text{ft}
\]
The maximum height \(H\) over which the water can be siphoned without cavitation occurring is approximately \(33.3 \, \text{ft}\). 
\newline
The maximum height above the tank at which water can be siphoned out is \( 33.3 - 12 = 21.3 \, \text{ft}\).

\newpage
\section*{Problem \# 6}

Water flows through the horizontal branching pipe shown in the Fig at a rate of \(10 \, \text{ft}^3/\text{s}\). If viscous effects are negligible, determine the water speed at section (2), the pressure at section (3), and the flowrate at section (4).

\subsection*{Solution}
We assume change in potential energy is negligible. \newline
We start by applying Bernoulli's equation between points 1 and 2:
\[
P_1 + \frac{1}{2} \rho v_1^2 = P_2 + \frac{1}{2} \rho v_2^2
\]
Given that \(P_1 = 10 \, \text{psi}\), \(P_2 = 5 \, \text{psi}\), \(A_1 = 1 \, \text{ft}^2\), and \(A_2 = 0.07 \, \text{ft}^2\).
\newline
The velocity at section 1 is \(v_1 = \dfrac{Q}{A_1} = \dfrac{10 \, \text{ft}^3/\text{s}}{1 \, \text{ft}^2} = 10 \, \text{ft/s}\).
\newline
Now, we can solve for \(v_2\):
\[
v_2 = \sqrt{2 \left( \frac{P_1 - P_2}{\rho} + \frac{v_1^2}{2} \right)}
\]
\[
v_2 = \sqrt{2 \left( \frac{(10 \, \text{psi} - 5 \, \text{psi}) \times 144 \, \text{in}^2/\text{ft}^2}{62.4 \, \text{lb/ft}^3 \times 1/32.2 } + \frac{(10 \, \text{ft/s})^2}{2} \right)}
\]
\[
v_2 = 29.0357 ft/s
\]
Assuming water has a density \(\rho = 62.4 \, \text{lb/ft}^3\), we can plug in the numbers to find \(v_2\).
\newline\newline
Next, we apply Bernoulli's equation between points 1 and 3:
\[
P_1 + \frac{1}{2} \rho v_1^2 = P_3 + \frac{1}{2} \rho v_3^2
\]
Given that \(v_3 = 20 \, \text{ft/s}\) and \(A_3 = 0.2 \, \text{ft}^2\), we can solve for \(P_3\):
\[
P_3 = P_1 + \frac{1}{2} \rho \left( v_1^2 - v_3^2 \right)
\]
\[
P_3 = (10 \, \text{psi} \times 144 \, \text{lb}_f/\text{ft}^2) + \frac{1}{2} \times 62.4 \, \text{lb}_m/\text{ft}^3 \times 1/32.2 \times \left( (10 \, \text{ft/s})^2 - (20 \, \text{ft/s})^2 \right)
\]
\[
P_3 = 1149.316 \, \text{lb}_f/\text{ft}^2
\]
\[
P_3 = 7.98 \, \text{psi}
\]
Finally, for section 4, we use the conservation of mass to find the flow rate \( Q_4 \):
\[
Q = Q_2 + Q_3 + Q_4
\]
We can use \( Q_2 = A_2 \times v_2 \) and \( Q_3 = A_3 \times v_3 \) to solve for \( Q_4 \).
\[
Q_2 = 0.07 \, \text{ft}^2 \times 29.0357 \, \text{ft/s} = 2.0325 \, \text{ft}^3/\text{s}
\]
\[
Q_3 = 0.2 \, \text{ft}^2 \times 20 \, \text{ft/s} = 4 \, \text{ft}^3/\text{s}
\]
\[
Q_4 = 10 \, \text{ft}^3/\text{s} - 2.0325 \, \text{ft}^3/\text{s} - 4 \, \text{ft}^3/\text{s} = 3.9675 \, \text{ft}^3/\text{s}
\]

\newpage
\section*{Problem \# 7}
Water (assumed inviscid and incompressible) flows steadily with a speed of 10 ft/s from the large tank. Determine the depth, \( H \), of the layer of light liquid (specific weight = 50 lbf/ft\(^3\)) that covers the water in the tank.

\subsection*{Solution}

Using the general equation for energy conservation between two points in a fluid flow:
\[
P_1 + \frac{1}{2} \rho_1 V_1^2 + \rho_1 g h_1 = P_2 + \frac{1}{2} \rho_2 V_2^2 + \rho_2 g h_2
\]
Given that \( \rho_1 = 1.94 \, \text{slug/ft}^3 \) for water, \( \rho_2 = 1.94 \, \text{slug/ft}^3 \) for water, and \( \rho_{\text{light liquid}} = \frac{50}{32.2} \, \text{slug/ft}^3 \), we can plug these into the equation:

\[
0 + \frac{1}{2} \times 1.94 \times 0^2 + 1.94 \times 32.2 \times 4 + \frac{50}{32.2} \times 32.2 \times H = 0 + \frac{1}{2} \times 1.94 \times 10^2 + 1.94 \times 32.2 \times 5
\]
Simplifying, we get:
\[
1.94 \times 32.2 \times 4 + 50 \times H = \frac{1}{2} \times 1.94 \times 100 + 1.94 \times 32.2 \times 5
\]
Further simplification leads to:
\[
1.94 \times 32.2 \times 4 + 50 \times H = 97 \times 1 + 1.94 \times 32.2 \times 5
\]
\[
H = 3.188 \, \text{ft}
\]
Therefore, the depth \( H \) of the layer of light liquid that covers the water in the tank is approximately 3.188 ft.

\newpage

\end{document}
